% Narragansett (Williams 1643) — Source Workflow
% Issue #1362
% Status: Research complete, implementation pending

\section{Narragansett: Williams 1643}
\label{sec:1362}

\subsection{Source Description}

\textbf{Recorder:} Roger Williams (1603--1683), founder of Rhode Island Colony.
\textbf{Source:} \emph{A Key into the Language of America} (1643), London: Gregory Dexter.
\textbf{Language:} Narragansett [xnt] (Southern New England Algonquian).
\textbf{Corpus:} 2,134 records.
\textbf{Orthography:} English-based ad-hoc spelling, seventeenth-century Early
Modern English conventions. Williams marks stress with diacritics (á primary,
à low/falling, â rising-falling) --- a major advantage over other colonial
sources such as Eliot's Massachusett Bible (1663), which uses no stress
marking.

\subsubsection{Recorder Biography}

Roger Williams (1603--1683) was a Puritan minister, theologian, and founder
of the Colony of Rhode Island and Providence Plantations. A graduate of
Pembroke College, Cambridge (B.A. 1627), Williams was a skilled linguist who
reportedly mastered Latin, Greek, Hebrew, Dutch, French, and several
Algonquian languages. He arrived in Massachusetts Bay in 1631 and began
intensive work with the Narragansett people after his banishment from the
Massachusetts Bay Colony in 1635. His \emph{A Key into the Language of
America} (1643) was the first extended study of a Native American language
published in English. The work takes the form of a phrase book organized by
topic (greetings, family, hunting, etc.), with each Narragansett phrase
accompanied by an English translation and occasional commentary.

\subsubsection{Orthography System}

Williams used an English-based ad-hoc spelling system with no standardized
orthography. He applied English sound-spelling conventions (seventeenth-century
Early Modern English) to Narragansett, an Algonquian language with a radically
different phoneme inventory. The system was designed to be usable by English
colonists, not by Algonquianists --- it systematically collapses phonemic
distinctions that English lacks.

\subsubsection{Stress Marking}

Williams' orthography is \textbf{stress-marked}, unlike Eliot's Massachusett
--- this is a major advantage. Eliot's Massachusett Bible (1663) uses no
accent marks, making stress reconstruction for Massachusett far more difficult.
Williams' accents provide critical data about Narragansett prosody.

\begin{table}[h]
\centering
\begin{tabular}{ll}
\toprule
Mark & Meaning \\
\midrule
á & Primary stress (or high pitch --- debated) \\
à & Low/falling tone (secondary stress rarely written) \\
â & Rising-falling tone \\
\bottomrule
\end{tabular}
\caption{Stress marking in Williams' orthography}
\label{tab:1362-stress}
\end{table}

Aubin's proposed stress rule (from personal correspondence) posits that all
long vowels are stressed and all alternate even-numbered short vowels are
stressed. However, this rule requires correct vowel length classification,
which has not been achieved for extinct SNEA languages. The problem is
circular: vowel length is needed to assign stress, but stress is one of the
cues for vowel length. Costa (2007) discusses the prosodic structure of SNEA
languages and notes that Narragansett stress patterns appear similar to Loup A
and Massachusett, with a preference for stress on long syllables and
alternating rhythm on short syllables --- but the documentary record is too
thin to confirm this definitively.

\subsection{Vowel Inventory}

\subsubsection{Simple Vowels}

\begin{table}[h]
\centering
\begin{tabular}{llll}
\toprule
Williams & IPA & Narragansett Phoneme & Notes \\
\midrule
a & /ə/, /ɔː/, /aː/ & PA *a *eː *aː & Three-way ambiguity \\
ah & /aː/, /ɔː/ & PA *aː & Long variants \\
e & /iː/, /ɛ/, /ə/, ∅ & PA *iː *e *a & Four-way ambiguity; silent word-final \\
ê & /iː/ & PA *iː & Circumflex = long i \\
i & /ə/, /iː/, /ɪ/ & PA *a *iː *i & Often = schwa in unstressed \\
o & /ə/, /aː/, /uː/; after w: /aː/, /ɔː/ & PA *a *aː *oː & Broad range \\
oo, ô & /uː/ & PA *oː & Consistent long u \\
u & /ə/, /a/; word-initially possibly /w/ & PA *a & Ambiguous schwa-ish \\
ie & /iː/ & PA *iː & \\
ee & /iː/ & PA *iː & Consistent \\
ei & /iː/, /ə/ & PA *iː *a & \\
ih & /ə/ & PA *a & \\
oh & /ə/, /oʊ/ & PA *a & \\
eu & /juː/ & --- & Rare, from English loan phonology \\
ea & /iː/, /ɛə/, /aː/ & PA *iː *aː & \\
\bottomrule
\end{tabular}
\caption{Simple vowel mappings in Williams' orthography}
\label{tab:1362-vowels}
\end{table}

\subsubsection{Nasalized Vowels}

A critical feature of SNEA phonology: Proto-Algonquian *aː and *oː developed
nasalized reflexes in Narragansett before certain consonants. Williams does
not mark nasality directly; he uses nasal coda consonant digraphs to imply it.

\begin{table}[h]
\centering
\begin{tabular}{ll}
\toprule
Williams & IPA \\
\midrule
an, aum, aun & /ã/ (PA *aː → ą in some environments) \\
om, on & /õ/ (PA *oː → ǫ) \\
ook & /oːk/, /oːkʷ/ (may be nasalized in context) \\
\bottomrule
\end{tabular}
\caption{Nasalized vowel patterns}
\label{tab:1362-nasal}
\end{table}

The nasalization is phonemic: \textit{n now} (yes) vs. \textit{núnow} (no)
--- only nasality distinguishes them, and Williams collapses this distinction
orthographically in related forms.

\subsubsection{Diphthong-Like Sequences}

\begin{table}[h]
\centering
\begin{tabular}{lll}
\toprule
Williams & IPA & Notes \\
\midrule
au & /ɔː/, /aʊ/ & \\
aw & /ɔː/, /ɒ/ & \\
ai & /aː/, /eɪ/ & \\
aü & /aː.uː/ & Rare; hiatal /aː/ + /uː/, not a true diphthong \\
\bottomrule
\end{tabular}
\caption{Diphthong-like sequences}
\label{tab:1362-diphthongs}
\end{table}

\subsection{Consonant Inventory}

SNEA languages neutralized the PA voicing distinction in stops. Williams'
English ear interpreted this as a gradient between voiced and voiceless,
writing both \textit{p} and \textit{b} (though \textit{b} is rare),
\textit{t} and \textit{d}, \textit{k} and \textit{g} (nearly absent). In
almost all cases, \textit{p}, \textit{t}, \textit{k} represent an unaspirated
stop --- between English /b/ and /p/, /d/ and /t/, /ɡ/ and /k/.

\subsubsection{Stops}

\begin{table}[h]
\centering
\begin{tabular}{lll}
\toprule
Williams & IPA & Notes \\
\midrule
p, pp & /b/ ~ /p/ [p] & Unaspirated lenis stop; alternates with \textit{b} \\
t & /d/ ~ /t/ ~ /tʲ/ & Before /i/ may palatalize to /tʃ/ \\
tt & /t/ ~ /d/ ~ /tʲ/ & Geminate, also palatalized in palatal contexts \\
k, kk & /k/ ~ /kʷ/ ~ /x/ & Labialized before /o/ /u/; guttural before /a/ \\
c, cc & /k/ ~ /kʷ/ & Before a/o/u = /k/; before e/i = /s/ (English soft-c) \\
\bottomrule
\end{tabular}
\caption{Stop consonants}
\label{tab:1362-stops}
\end{table}

Cowan (1973) demonstrated that PA *-k- (simple velar) and *-hk- (preaspirated
velar) merged phonetically in Narragansett as /hk/ medially, but can sometimes
be distinguished through the behavior of preceding vowels. PA *-k- shortens a
preceding long vowel; PA *-hk- does not. This provides an indirect method for
reconstructing preaspiration from vowel length patterns --- even though
Williams never writes the /h/.

\subsubsection{Affricates}

\begin{table}[h]
\centering
\begin{tabular}{lll}
\toprule
Williams & IPA & Notes \\
\midrule
ch & /tʃ/ & As in English ``chair'' --- may also represent /t/ palatalized \\
tch & /tʃ/ & As in English ``itch'' \\
ddt, dt & /d/ ~ /t/ ~ /tʲ/ & Complex ch-t sound; among the most ambiguous \\
\bottomrule
\end{tabular}
\caption{Affricate consonants}
\label{tab:1362-affricates}
\end{table}

\subsubsection{Nasals}

\begin{table}[h]
\centering
\begin{tabular}{lll}
\toprule
Williams & IPA & Notes \\
\midrule
m, mm & /m/ & \\
n, nn & /n/ & Before consonant can be syllabic /n̩/ (→ /nə/ for English speakers) \\
n (syllabic) & /n̩/ → /nə/ & Unwritten schwa before n in consonant clusters \\
\bottomrule
\end{tabular}
\caption{Nasal consonants}
\label{tab:1362-nasals}
\end{table}

\subsubsection{Fricatives}

\begin{table}[h]
\centering
\begin{tabular}{lll}
\toprule
Williams & IPA & Notes \\
\midrule
s, ss & /s/ & Single \textit{s} after vowel; geminate \textit{ss} also /s/ \\
sh & /ʃ/ word-initially; /s/ before consonant & Context-dependent ambiguity \\
z & /s/ (not /z/) & English-style voiced s for fortis /s/ \\
sc, shk, sk, shq & /sk/ & Cluster spellings \\
sq & /skʷ/ ~ /x/ & Labialized or guttural depending on following vowel \\
\bottomrule
\end{tabular}
\caption{Fricative consonants}
\label{tab:1362-fricatives}
\end{table}

Goddard \& Bragdon (1988) note that in Massachusett, the cognate system shows
*/s/ leniting to /h/ in intervocallic position --- a sound change that may
also operate in Narragansett but is invisible in Williams' orthography (since
he never writes /h/).

\subsubsection{Glides}

\begin{table}[h]
\centering
\begin{tabular}{lll}
\toprule
Williams & IPA & Notes \\
\midrule
w, ww & /w/ ~ /j/ before e/i & Before e/i can be /j/ (yes) \\
q (before vowel) & /kʷ/ & As in English ``queen'' --- consistent \\
q (before w) & /k/ ~ /kʷ/ ~ /x/ & Context-dependent; may represent guttural \\
\bottomrule
\end{tabular}
\caption{Glide consonants}
\label{tab:1362-glides}
\end{table}

\subsubsection{The Silent /h/ --- Preaspiration}

\textbf{English speakers routinely fail to hear preaspirated /h/ before stops.}
This is the single most important finding for this source. Williams almost
never wrote /h/ where modern reconstructions put preaspiration:

\begin{itemize}
  \item In Narragansett (and other SNEA languages), /h/ commonly precedes /p/,
    /t/, /k/ in word-medial position (e.g., /hkaː/, /hp/, /ht/)
  \item English speakers perceive these as single stops /k/, /p/, /t/ ---
    preaspiration is not a phonemic distinction in English
  \item Williams wrote \textit{k} where the actual pronunciation was /hk/
    (preaspirated k)
  \item \textbf{Confirmed by \textbackslash ph records:} The 3 records with
    \textbackslash ph fields show missing /h/ in Williams
\end{itemize}

\begin{table}[h]
\centering
\begin{tabular}{llll}
\toprule
ID & Williams & \textbackslash ph (phonetic) & Pattern \\
\midrule
3834 & Aaunchemókaw & ąːčimoːhkaːw & \textit{k} → hk \\
4286 & Anakáusichick & anoːhkaːsiːčiːk & \textit{k} → hk \\
4285 & Anakáusu & anoːhkaːsəw & \textit{k} → hk \\
\bottomrule
\end{tabular}
\caption{Evidence from the 3 \textbackslash ph entries in the database}
\label{tab:1362-ph}
\end{table}

All three show the same pattern: Williams writes \textit{k} where the
reconstruction has /hk/. The /h/ is invisible to English ears.

\subsection{Problems Encountered}

\subsubsection{Vowel Quality Ambiguity}

Williams' vowels \textit{a}, \textit{e}, \textit{i}, \textit{o}, \textit{u}
each map to 2--4 possible phonemes. The vowel \textit{a} alone can represent
/ə/, /ɔː/, or /aː/. This extreme ambiguity means that any single Williams
spelling is consistent with multiple phonological interpretations.

\subsubsection{Nasalization Implied but Not Marked}

Sequences like \textit{an}, \textit{aun}, \textit{om} may indicate nasal
vowels /ã/, /õ/ or oral vowel + consonantal \textit{n}. Disambiguating
requires:
\begin{enumerate}
  \item Checking the cognate in Massachusett (Eliot marks nasal /ô/ explicitly)
  \item Checking PA reconstruction (do expected reflexes predict nasalization
    in this environment?)
  \item Checking morphological structure (is the \textit{n} part of a suffix,
    or a nasalization marker?)
\end{enumerate}

\subsubsection{Stop Voicing Neutralized}

\textit{p}, \textit{t}, \textit{k} represent unaspirated stops that English
speakers perceived variably as voiced or voiceless. Williams writes both
\textit{p} and \textit{b}, \textit{t} and \textit{d}, \textit{k} and
\textit{g} in free variation.

\subsubsection{English Spelling Conventions}

English digraphs (\textit{ea}, \textit{oo}, \textit{ch}, \textit{sh}) follow
English values, but inconsistently due to 17th-century English pronunciation
being different from modern English. For example, \textit{ea} in 1643 was
/ɛː/, not modern /iː/.

\subsubsection{Silent e}

Word-final \textit{e} is often silent, but not always. This affects at least
15--20\% of entries. Comparison with Massachusett cognates (Eliot) is the
primary disambiguation strategy.

\subsubsection{Preaspiration /h/ Invisible}

Only 3 of 2,134 records have phonetic (\textbackslash ph) fields. All three
show the same pattern: Williams writes \textit{k} where the reconstruction has
/hk/. If this pattern is representative, >50\% of Williams' stops may be
missing a preceding /h/.

\subsubsection{\textit{sh} Before Consonants}

Williams writes \textit{sh} before consonants where Massachusett cognates
suggest plain /s/: Williams \textit{shp-} corresponds to Massachusett
\textit{sp-}. Costa (2007) notes that this palatalization before consonants is
also attested in Loup A records, suggesting it is a genuine SNEA dialect
feature, not a Williams idiosyncrasy.

\subsubsection{The \textit{an}/\textit{aun} Ambiguity for /ã/}

The most consequential ambiguity: Williams writes both \textit{an} and
\textit{aun} for /ã/, but also writes \textit{an} for plain /an/ (oral a + n).
Disambiguating requires comparative evidence from Massachusett (Eliot marks
nasal /ô/ explicitly), PA reconstruction, and morphological analysis.

\subsection{Research Approach}

\subsubsection{Recommended Methodology}

\textbf{DO NOT USE REGEX} --- the ambiguities are fundamentally
context-dependent. Recommended approach:

\begin{enumerate}
  \item \textbf{Finite State Machine (FSM):} Encode known spelling-to-phoneme
    rules as weighted state transitions, with fallback probabilities for
    ambiguous mappings.
  \item \textbf{ML classifier:} Train on the 3 \textbackslash ph seed examples
    (plus indirect evidence from comparative Algonquian) to predict
    preaspiration placement and vowel quality.
  \item \textbf{Morphological analysis:} Use known Narragansett morphology
    (prefixes, stems, suffixes from Aubin/Hagenau) to segment Williams'
    spellings before applying conversion.
  \item \textbf{Comparative evidence:} Use Proto-Algonquian reconstructions
    and sister languages (Massachusett, Mohegan-Pequot) as priors for
    ambiguous mappings.
\end{enumerate}

\subsubsection{Annotation Plan}

The 3 \textbackslash ph entries alone are insufficient for ML training.
Recommended annotation plan:

\begin{enumerate}
  \item Hand-annotate ~50 representative entries with conversions
  \item Use the annotations as training data
  \item Evaluate on the remaining entries with spot-check verification
  \item Cross-validate against Massachusett cognates (Goddard \& Bragdon 1988)
    for at least 200 entries to calibrate the Narragansett-specific isoglosses
\end{enumerate}

\subsubsection{Cross-Linguistic Cognate Evidence}

Massachusett, the closest sister language to Narragansett, was documented by
John Eliot about 20 years after Williams. Eliot used a different orthographic
convention --- less English-phonetic, more systematic (he devised it for Bible
translation).

\begin{table}[h]
\centering
\begin{tabular}{llll}
\toprule
English & Williams Narragansett & Eliot Massachusett & PA reconstruction \\
\midrule
dog & ayim & arum & *aɬemwa \\
man & weenawash & wénawaš & \\
water & nipp & nippe & *nepyi \\
fire & apeisuwau & apeisuwâ & \\
house & wétu & wétu & *wiːkiwaːmi \\
I (say) & neen & neen & *niːra \\
\bottomrule
\end{tabular}
\caption{Cognate pairs for calibration}
\label{tab:1362-cognates}
\end{table}

\subsubsection{Open Questions for Revitalization}

\begin{enumerate}
  \item \textbf{Vowel length:} Was Narragansett vowel length phonemic (as in
    Massachusett)? Williams' orthography suggests yes (áh vs. a for long vs.
    short /a/), but the marking is inconsistent.
  \item \textbf{Preaspiration:} How pervasive was /h/ before stops? The 3
    \textbackslash ph entries suggest it was common.
  \item \textbf{Nasalization phonemicity:} Was nasalization phonemic (as in
    Massachusett) or allophonic? Costa (2007) suggests phonemic in at least
    some SNEA dialects.
  \item \textbf{Final vowel status:} Did Narragansett drop word-final vowels
    (like Mahican) or retain them (like Massachusett)? Williams' variable
    practice suggests either retention with devoicing or transition state.
\end{enumerate}

\subsection{Conversion Workflow}
\texttt{[Pending implementation --- see Issue \#1362]}

\subsection{Linguist Feedback}
\texttt{[Template --- content pending review]}

\subsection{Related Work: Colonial Recorder Perception Problem}

The core challenge of this research card --- English-speaking recorders
misperceiving and mis-transcribing Algonquian phonemes --- is a known
phenomenon documented across Eastern North America. Two researchers' work is
directly relevant:

\textbf{Dr. Keith Cunningham} (Linguist, Nanticoke Indian Tribe). His doctoral
dissertation \emph{A Phonological Analysis of Nanticoke With Practical
Applications for Language Revitalization} (Georgetown University) analyzes
three 18th-century Nanticoke word lists (Heckewelder 1785) using the same
methodology: historical phonology from colonial recorders, with the same
challenges of multilingual informants and orthographic complexity. His 2024
Algonquian Conference presentation \emph{A Revised Phonological Analysis of
the Heckewelder Vocabulary of Nanticoke} directly parallels the Williams
analysis --- both involve English-speaking recorders working without
standardized orthography, both require PA comparative calibration, and both
feed into practical language revitalization. Cunningham's work demonstrates
that the colonial recorder perception problem extends beyond SNEA into the
broader Eastern Algonquian family (Nanticoke, Maryland Algonquian).

\textbf{Dr. Craig Kopris} (Independent scholar). His paper \emph{Les sons
wyandots perçus par des oreilles étrangères} (Wyandot sounds perceived by
foreign ears, \emph{Recherches amérindiennes au Québec}, 2014) is the exact
framing of the Williams problem: how European recorders systematically
misperceived and mis-transcribed indigenous phonemes. His \emph{Wyandot
Phonology: Recovering the Sound System of an Extinct Language} applies the
same recovery methodology to an unrelated language family (Iroquoian),
demonstrating that the foreign-ear perception problem is universal across
colonial documentation of North American languages. His Dictionary of Wyandot
project (NEH award FN-255576-17) shows the long-term application of these
methods to language revitalization.

The Williams Narragansett corpus (2,134 records, only 3 with phonetic
transcription) represents one of the most extreme cases of this problem ---
but the methodological framework for addressing it is established in the
parallel work of Cunningham (Algonquian) and Kopris (Iroquoian).

\subsection{Bibliography}

\begin{enumerate}

\item Bloomfield, Leonard. (1925). ``On the Sound-System of Central
Algonquian.'' \emph{Language} 1(4): 130--156.

\item Bloomfield, Leonard. (1946). ``Algonquian.'' In Harry Hoijer et al.,
\emph{Linguistic Structures of Native America}, pp. 85--129. Viking Fund
Publications in Anthropology 6. New York: Wenner-Gren Foundation.

\item Bragdon, Kathleen J. (1999). \emph{Native People of Southern New
England, 1500--1650}. Norman: University of Oklahoma Press.

\item Bragdon, Kathleen J. (2009). \emph{Native Writings in Massachusett,
Part II: Linguistic Analysis}. Memoirs of the American Philosophical Society
263. Philadelphia: American Philosophical Society.

\item Costa, David J. (2007). ``The Dialectology of Southern New England
Algonquian.'' In H.C. Wolfart, ed., \emph{Papers of the 38th Algonquian
Conference}, pp. 81--127. Winnipeg: University of Manitoba.

\item Cowan, William. (1973). ``The Development of PA *-k- and *-hk- in
Narragansett.'' \emph{International Journal of American Linguistics} 39(2):
90--93.

\item Cowan, William. (1975). ``PA *-θ- and *-t- in Narragansett.''
\emph{International Journal of American Linguistics} 41(2): 178--181.

\item Goddard, Ives. (1978). ``Eastern Algonquian Languages.'' In Bruce G.
Trigger, ed., \emph{Handbook of North American Indians, Vol. 15: Northeast},
pp. 70--77. Washington: Smithsonian Institution.

\item Goddard, Ives. (1996). ``The Description of the Native Languages of
North America Before Boas.'' In Ives Goddard, ed., \emph{Handbook of North
American Indians, Vol. 17: Languages}, pp. 17--42. Washington: Smithsonian
Institution.

\item Goddard, Ives, and Kathleen J. Bragdon. (1988). \emph{Native Writings
in Massachusett.} 2 vols. Memoirs of the American Philosophical Society 185.
Philadelphia: American Philosophical Society.

\item Hagenau, Walter. (1962). \emph{A Morphological Study of Narragansett.}
M.A. thesis, University of North Dakota.

\item O'Brien, Frank Waobu. \emph{Appendix: Guide to Historical Spellings
and Sounds in New England Algonquian Languages.} Aquidneck Indian Council.

\item Pentland, David H. (1979). \emph{Algonquian Historical Phonology.}
Ph.D. dissertation, University of Toronto.

\item Siebert, Frank T., Jr. (1975). ``Resurrecting Virginia Algonquian from
the Dead: The Reconstituted and Historical Phonology of Powhatan.'' In James
M. Crawford, ed., \emph{Studies in Southeastern Indian Languages}, pp.
285--453. Athens: University of Georgia Press.

\item Trumbull, James Hammond. (1870). ``Notes on Forty Versions of the
Lord's Prayer in Algonkin Languages.'' \emph{Transactions of the American
Philological Association} (1869--1896): 1--20.

\item Williams, Roger. (1643). \emph{A Key into the Language of America.}
London: Gregory Dexter.

\end{enumerate}
